% Exported from Math++. Compile with pdflatex.
\documentclass[11pt,a4paper]{article}
\usepackage[T1]{fontenc}
\usepackage[utf8]{inputenc}
\usepackage{lmodern}
\usepackage{amsmath,amssymb,amsthm,mathtools}
\usepackage{graphicx,booktabs,xcolor}
\usepackage[hidelinks]{hyperref}
\usepackage[margin=25mm]{geometry}
\theoremstyle{plain}
\newtheorem{theorem}{Theorem}
\newtheorem{lemma}[theorem]{Lemma}
\newtheorem{corollary}[theorem]{Corollary}
\newtheorem{proposition}[theorem]{Proposition}
\newtheorem{conjecture}[theorem]{Conjecture}
\theoremstyle{definition}
\newtheorem{definition}[theorem]{Definition}
\newtheorem{example}[theorem]{Example}
\theoremstyle{remark}
\newtheorem*{remark}{Remark}
\setlength{\parindent}{0pt}\setlength{\parskip}{0.6em}
\title{Multivariable and vector calculus}
\date{}
\begin{document}
\maketitle
\section*{Partial derivatives and the gradient}

\begin{equation*}
  f(x, y) \coloneqq  x^2 y + y^3
\end{equation*}
\begin{equation*}
  \frac{\partial^2 f}{\partial x \partial y} = 2 x
\end{equation*}
\begin{equation*}
  \operatorname{grad}_{(x,y)=(1,2)} f = \left(4, 13\right)
\end{equation*}
\begin{equation*}
  \operatorname{Directional\ derivative}_{(x,y)=(1,2)}(f, (3, 4)) = \frac{64}{5}
\end{equation*}
\section*{Tangent planes and linear approximation}

\begin{equation*}
  \operatorname{Tangent\ line}_{(x,y)=(1,2)} (x^2 + y^2) = z = 2 x + 4 y - 5
\end{equation*}
\begin{equation*}
  \operatorname{Tangent\ plane}_{(x,y,z)=(1,2,3)} (x^2 + y^2 + z^2 = 14) = 2 x + 4 y + 6 z - 28 = 0
\end{equation*}
\begin{equation*}
  \operatorname{Linearisation}_{(x,y)=(3,4)} \sqrt{x^2 + y^2} = L\left(x, y\right) = \frac{3 x}{5} + \frac{4 y}{5}
\end{equation*}
\section*{Extrema}

The second-derivative test classifies each critical point:

\begin{align*}
  \operatorname{Second-derivative\ test} (x^3 + y^3 - 3xy) = \begin{array}{ll} \text{saddle point at (0, 0)} & 0 \\ \text{local minimum at (1, 1)} & -1 \end{array}
\end{align*}
Lagrange multipliers on the circle $x^2 + y^2 = 1$:

\begin{align*}
  \operatorname{Lagrange\ multipliers}(x y, x^2 + y^2 = 1) = \begin{array}{ll} \text{f(-sqrt(2)/2, sqrt(2)/2)} & -\frac{1}{2} \\ \text{f(sqrt(2)/2, -sqrt(2)/2)} & -\frac{1}{2} \\ \text{f(-sqrt(2)/2, -sqrt(2)/2)} & \frac{1}{2} \\ \text{f(sqrt(2)/2, sqrt(2)/2)} & \frac{1}{2} \\ \text{largest value} & \frac{1}{2} \\ \text{smallest value} & -\frac{1}{2} \end{array}
\end{align*}
\section*{Integrals in other coordinates}

The Jacobian ($r$, $r$ and $\rho^2 \sin\phi$) is put in for you.

\begin{equation*}
  \operatorname{polar}(x^2 + y^2, \theta \in [0, 2\pi], r \in [0, 2]) = 8 \pi
\end{equation*}
\begin{equation*}
  \operatorname{cylindrical}(z, \theta \in [0, 2\pi], r \in [0, 1], z \in [0, 2]) = 2 \pi
\end{equation*}
\begin{equation*}
  \operatorname{spherical}(1, \theta \in [0, 2\pi], \phi \in [0, \pi], \rho \in [0, 1]) = \frac{4 \pi}{3}
\end{equation*}
\begin{equation*}
  \operatorname{Change\ of\ variables}(x + y, (x, y) = (u + v, u - v), u \in [0, 1], v \in [0, 1]) = 2
\end{equation*}
\section*{Vector fields}

\begin{equation*}
  \nabla \cdot (x y, y z, z x) = x + y + z
\end{equation*}
\begin{equation*}
  \nabla \times (x y, y z, z x) = \left(-y, -z, -x\right)
\end{equation*}
\begin{equation*}
  \nabla^2 (x^2 + y^2 + z^2) = 6
\end{equation*}
\begin{equation*}
  \operatorname{Line\ integral}((-y, x), (\cos t, \sin t), t \in [0, 2\pi]) = 2 \pi
\end{equation*}
\begin{equation*}
  \operatorname{Flux}((x, y, z), (\sin u \cos v, \sin u \sin v, \cos u), u \in [0, \pi], v \in [0, 2\pi]) = 4 \pi
\end{equation*}
\section*{The integral theorems, both sides worked out}

\begin{align*}
  \operatorname{Green's\ theorem}((-y, x), x \in [0, 1], y \in [0, 1]) = \begin{array}{ll} \text{\oint  P dx + Q dy} & 2 \\ \text{\iint  (\partial Q/\partial x ? \partial P/\partial y) dA} & 2 \\ \text{the two sides agree} & \text{yes} \end{array}
\end{align*}
\begin{align*}
  \operatorname{Stokes'\ theorem}((-y, x, 0), (\sin u \cos v, \sin u \sin v, \cos u), u \in [0, \pi/2], v \in [0, 2\pi]) = \begin{array}{ll} \text{\oint  F \cdot  dr} & 2 \pi \\ \text{\iint  (\nabla  \times  F) \cdot  dS} & 2 \pi \\ \text{the two sides agree} & \text{yes} \end{array}
\end{align*}
\begin{align*}
  \operatorname{Divergence\ theorem}((x, y, z), \rho \in [0, 1], \phi \in [0, \pi], \theta \in [0, 2\pi]) = \begin{array}{ll} \text{\oiint  F \cdot  dS} & 4 \pi \\ \text{\iiint  \nabla  \cdot  F dV} & 4 \pi \\ \text{the two sides agree} & \text{yes} \end{array}
\end{align*}
\section*{Implicit differentiation and related rates}

\begin{equation*}
  \operatorname{Implicit\ differentiation}_x (x^2 + y^2 = 25) = \frac{d}{d x}\left(y\right) = -\frac{x}{y}
\end{equation*}
A ladder of length 5: the foot slides away at 2 m/s when it is 3 m from the wall.

\begin{equation*}
  \operatorname{rates}_t (x^2 + y^2 = 25, x = 3, y = 4, dx = 2) = \frac{d}{d t}\left(y\right) = -\frac{3}{2}
\end{equation*}
\section*{Laplace transforms}

\begin{equation*}
  \mathcal{L} t^2 e^{3t} = \frac{2}{\left(s - 3\right)^{3}}
\end{equation*}
\begin{equation*}
  \mathcal{L}^{-1} \frac{s + 1}{s^2 + 2s + 5} = e^{-t} \cos\left(2 t\right)
\end{equation*}
Solve $y'' - 3y' + 2y = 0$ with $y(0) = 0$ and $y'(0) = 1$ (coefficients from $y$ up):

\begin{equation*}
  \operatorname{Solve\ by\ Laplace\ transform}((2, -3, 1), 0, (0, 1)) = y\left(t\right) = -e^{t} + e^{2 t}
\end{equation*}
\end{document}
